% ====================================================================
% Chapter 5: Evaluation Systems and Benchmarks (English expanded version)
% ====================================================================

\section{Evaluation Systems and Benchmarks}
\label{ch:evaluation}

\subsection{The Evaluation Challenge in Agent Self-Evolution}
\label{sec:eval-challenge}

The central tension in agent self-evolution research lies between a system's claim of improvement and our ability to verify that improvement is genuine, robust, and transferable. Evaluation is not merely the experimental section of a paper---it serves as the \emph{selection pressure} that shapes the direction of evolution and the \emph{safety gate} that prevents harmful mutations from propagating. If the evaluation function is misaligned, incomplete, or gameable, then self-evolution becomes self-deception.

This chapter addresses the evaluation problem from three perspectives. First, we propose a five-layer evaluation hierarchy that spans fully automated programmatic verification at the bottom to human-in-the-loop safety gates at the top. Second, we survey the major benchmarks used across the agent evolution literature, analyzing their strengths, limitations, and contamination risks. Third, we introduce agent-specific evaluation metrics that go beyond single-task accuracy to capture iterative improvement, transfer, cost efficiency, and safety. Throughout, we emphasize that evaluation infrastructure is itself evolving: as agents become more capable, evaluators must become more sophisticated, or risk being outpaced by the systems they are meant to assess.

The evaluation challenge is particularly acute for self-evolving systems because of four properties that distinguish them from conventional machine learning. First, the improvement target is moving: an agent that modifies its own code, prompts, or memory may change its behavior in ways that the benchmark did not anticipate. Second, the evaluation is part of the optimization loop: agents can learn to exploit evaluator weaknesses, producing reward hacking, length bias, or score manipulation that inflates metrics without genuine capability gain. Third, the scope of evaluation must expand beyond task completion to include process quality, cost, safety, and transfer---dimensions that are harder to quantify than accuracy. Fourth, reproducibility is threatened by non-determinism in LLM generation, API version changes, and the interaction between stochastic sampling and evolving agent state.

These properties mean that a single benchmark score is never sufficient evidence for self-evolution. A credible evaluation must include multi-layer verification, cross-benchmark generalization, cost analysis, safety auditing, and ablation studies that isolate the contribution of the self-improvement mechanism from confounds such as prompt engineering, model selection, and data contamination.


\subsection{A Five-Layer Evaluation Hierarchy}
\label{sec:eval-hierarchy}

We propose organizing evaluation methods into five layers, from fully automated to human-supervised. Each layer adds reliability at the cost of increased effort and latency. A well-designed evaluation system uses all five layers in combination, with lower layers providing rapid feedback during development and higher layers serving as pre-deployment gates.

% ====================================================================
% Figure 5.1: Evaluation Benchmark Layered System
% ====================================================================
\begin{figure}[htbp]
\centering
\begin{tikzpicture}[
    layer/.style={rectangle, rounded corners=6pt, draw=#1!70!black, fill=#1!12,
        minimum width=12cm, minimum height=1.4cm, align=center, font=\small},
    bench/.style={font=\scriptsize, color=gray!60},
]
    % Layers from bottom to top
    \node[layer=cat1] (l1) at (0,0) {\textbf{L1: Programmatic Verification} --- Unit Tests, Type Checking, Linting, Fuzzing, Sandbox Execution};
    \node[layer=cat2] (l2) at (0,2) {\textbf{L2: Code \& Reasoning Benchmarks} --- SWE-Bench, HumanEval, GSM8K, MATH, MMLU, GPQA};
    \node[layer=cat3] (l3) at (0,4) {\textbf{L3: Environment Interaction Benchmarks} --- WebArena, ALFWorld, Voyager/Minecraft, AppWorld};
    \node[layer=cat4] (l4) at (0,6) {\textbf{L4: Process Metrics} --- Iterative Gain Curves, Archive Diversity, Transferability, Rollback Rate};
    \node[layer=cat5] (l5) at (0,8) {\textbf{L5: Safety \& Cost Gates} --- Permission Checks, Budget Limits, Regression Tests, Human Audit};

    % Side annotations
    \node[font=\scriptsize\bfseries, color=cat1!60!black, rotate=90] at (-7,0) {Automated};
    \node[font=\scriptsize\bfseries, color=cat5!60!black, rotate=90] at (-7,8) {Human};
    \draw[-{Stealth}, thick, dashed, color=gray!40] (-6.5,0) -- (-6.5,8);

\end{tikzpicture}
\caption{Five-layer evaluation hierarchy for agent self-evolution. From L1 (fully automated programmatic verification) to L5 (human-in-the-loop safety gates), evaluation cost and reliability increase at each layer.}
\label{fig:eval-layers}
\end{figure}


\subsubsection{L1: Programmatic Verification}
\label{sec:eval-L1}

The most reliable evaluation signals come from deterministic programmatic checks. These include unit tests, type checking, static analysis, linting, fuzzing, regression testing, and sandbox execution. Their key advantage is that they are fully automated, reproducible, and cannot be gamed by an LLM that understands the evaluation format---assuming the test suite itself is not accessible to the agent.

In the agent evolution literature, systems like DGM~\citep{dgm2025}, SICA~\citep{sica2025}, EvoMAC~\citep{evomac2024}, and AlphaEvolve~\citep{alphaevolve2025} achieve self-evolution precisely because candidate code can be automatically executed and scored. DGM evaluates agent variants by running them on SWE-Bench tasks and measuring the percentage of tests passed; SICA allows the coding agent to modify its own source code and validates changes against a held-out test suite; AlphaEvolve uses programmatic evaluators to score algorithm candidates on mathematical optimization problems.

The scope of programmatic verification should be expanded beyond code tasks. Data pipelines can verify schema compliance, distribution drift, and end-to-end consistency. Business processes can verify state machine constraints and permission policies. Web tasks can verify DOM state, backend responses, and user intent fulfillment. Programmatic verifiers do not need to understand full semantics, but they must cover the critical failure modes of the target domain.

A key challenge is test coverage: a test suite that passes may simply be incomplete. Self-evolving agents can learn to satisfy the visible tests while introducing subtle regressions in untested behaviors. This motivates mutation testing (injecting deliberate bugs to check whether tests catch them), fuzzing (generating random inputs to stress-test edge cases), and invariant checking (verifying properties that should hold across all executions).


\subsubsection{L2: Code and Reasoning Benchmarks}
\label{sec:eval-L2}

Code generation and mathematical reasoning benchmarks provide standardized, widely-recognized evaluation signals. Their primary advantage is comparability: results can be directly compared across papers and systems.

For code generation, HumanEval~\citep{chen2021codex} consists of 164 Python programming problems with unit tests, measuring functional correctness via pass@$k$ metrics. MBPP~\citep{austin2021mbpp} provides 974 crowd-sourced problems. Reflexion~\citep{reflexion2023} demonstrates self-improvement on HumanEval, achieving 91\% pass@1 with GPT-3.5-turbo, surpassing GPT-4's 80.1\% zero-shot. Self-Refine~\citep{selfrefine2023} shows approximately 20\% average improvement across 7 tasks including code optimization (+28\%), math reasoning (56\% to 67\%), and constrained generation (65\% to 78\%).

For mathematical reasoning, GSM8K~\citep{cobbe2021gsm8k} provides 8,500 grade-school math problems, MATH~\citep{hendrycks2021math} offers 12,500 competition-level problems, and GPQA~\citep{rein2024gpqa} presents expert-level questions. STaR~\citep{star2022} improves from 35\% to 74\% on GSM8K via self-taught reasoning. RISE~\citep{rise2024} improves Llama2-7B from approximately 14\% to 24\% and Mistral-7B from 38\% to 46\% through multi-turn RL introspection.

However, these benchmarks suffer from well-documented limitations. Many are static problem sets vulnerable to data contamination. They measure only final-answer correctness, not reasoning process quality, tool use, or multi-step interactions. Their granularity is small: passing a unit test on a 10-line function differs fundamentally from patching a real codebase.


\subsubsection{L3: Environment Interaction Benchmarks}
\label{sec:eval-L3}

Environment interaction benchmarks evaluate agents in simulated or real environments where success requires multi-step planning, tool use, error recovery, and state tracking.

SWE-Bench~\citep{jimenez2024swebench} requires systems to locate issues in real open-source repositories, modify code, run tests, and submit verifiable patches. DGM~\citep{dgm2025} improves from 20.0\% to 50.0\% on SWE-Bench Verified through archive-based self-evolution---a 2.5$\times$ relative improvement. DGM also improves from 14.2\% to 30.7\% on the Polyglot multi-language coding benchmark, demonstrating cross-language transfer.

WebArena~\citep{zhou2024webarena} creates realistic web-based tasks requiring website navigation, form filling, and information extraction. ALFWorld~\citep{shridhar2021alfworld} tests embodied agents in text-based household environments. Reflexion achieves 97\% on ALFWorld (up from 75\% baseline), demonstrating that language-based reflection is effective in multi-step interactive settings.

Voyager~\citep{voyager2023} uses unique evaluation metrics in Minecraft: unique items discovered, technology tree coverage, and cross-world skill transfer. This lifelong learning evaluation captures cumulative capability growth rather than single-test performance.


\subsubsection{L4: Process Metrics}
\label{sec:eval-L4}

Process metrics evaluate \emph{how} improvement occurs, distinguishing genuine self-evolution from one-time optimization.

\textbf{Iterative Improvement Curves.} We define $\Delta_t = V(z_t, \tau) - V(z_{t-1}, \tau)$ where $V(z_t, \tau)$ is the evaluation score of agent state $z_t$ on task distribution $\tau$. Key patterns include: monotonic improvement ($\Delta_t > 0$ for all $t$), plateau ($\Delta_t \approx 0$ after initial gains), volatile improvement (alternating positive and negative $\Delta_t$), and regression ($\Delta_t < 0$ indicating negative transfer).

\textbf{Transfer and Generalization.} ADAS~\citep{adas2024} demonstrates cross-model and cross-domain transfer: agent designs discovered on ARC transfer to GFootball, and designs searched with GPT-3.5 remain effective on GPT-4. Absolute Zero~\citep{absolutezero2025} shows that code self-play training confers mathematical reasoning ability---a striking cross-domain transfer result.

\textbf{Archive Diversity.} DGM's key insight is that maintaining a growing archive of agent variants---including currently suboptimal ones---enables open-ended exploration and avoids local optima. Archive diversity metrics quantify the breadth of the explored design space.

\textbf{Rollback Rate.} The frequency of rollbacks to previous versions directly measures evolution reliability. High rollback rates indicate insufficiently conservative mutation-selection processes.


\subsubsection{L5: Safety and Cost Gates}
\label{sec:eval-L5}

The highest layer evaluates whether improvement was achieved at acceptable cost and risk.

\textbf{Permission Verification.} Every candidate variant must undergo permission checks. Self-modifying agents like G\"{o}del Agent~\citep{godelagent2024} and SICA are particularly sensitive to this concern.

\textbf{Cost Accounting.} Key metrics: API calls per improvement iteration, token consumption per unit gain, wall-clock time to target performance, and human intervention frequency. Cost should be part of the fitness function.

\textbf{Regression Testing.} Every evolved variant must pass regression tests verifying core capability preservation. Self-evolution can produce negative transfer: optimizing for new tasks may degrade existing ones.

\textbf{Human Audit.} For high-stakes applications, human review of code changes, reasoning traces, and safety constraint maintenance provides the final safety net.


\subsection{Quantitative Benchmark Results Across Methods}
\label{sec:eval-quantitative}

This section consolidates quantitative results from the agent evolution literature.

\begin{table}[htbp]
\centering
\small
\caption{Code generation and software engineering benchmark results across self-improvement methods.}
\label{tab:eval-code}
\begin{tabular}{p{28mm}p{22mm}p{18mm}p{18mm}p{30mm}}
\toprule
\textbf{Method} & \textbf{Base Model} & \textbf{HumanEval} & \textbf{SWE-Bench} & \textbf{Key Feature} \\
\midrule
GPT-4 (zero-shot) & GPT-4 & 80.1\% & --- & Baseline \\
Reflexion & GPT-3.5-turbo & \textbf{91.0\%} & --- & Verbal RL + memory \\
Self-Refine & GPT-3.5/4 & +20\% avg & --- & Self-feedback loop \\
DGM (seed) & Claude/GPT & --- & 20.0\% & Initial agent \\
\textbf{DGM (evolved)} & Claude/GPT & --- & \textbf{50.0\%} & Archive evolution \\
AZR-Coder & Qwen-7B & SOTA (7B) & --- & Zero-data self-play \\
\bottomrule
\end{tabular}
\end{table}

Reflexion's 91\% pass@1 on HumanEval using GPT-3.5-turbo surpasses GPT-4 zero-shot (80.1\%) by 10.9 percentage points. The memory buffer---storing natural language descriptions of failures and corrections---acts as a ``verbal gradient'' guiding subsequent attempts. DGM's SWE-Bench improvement (20.0\% to 50.0\%, a 2.5$\times$ relative gain) demonstrates that agent-level self-modification can produce gains beyond prompt-level or memory-level self-improvement. Absolute Zero achieves SOTA at 7B scale using zero external training data, with code self-play generating both tasks and solutions verified by a code executor.

\begin{table}[htbp]
\centering
\small
\caption{Mathematical reasoning benchmark results across self-improvement methods.}
\label{tab:eval-math}
\begin{tabular}{p{28mm}p{22mm}p{18mm}p{18mm}p{30mm}}
\toprule
\textbf{Method} & \textbf{Base Model} & \textbf{GSM8K} & \textbf{MATH} & \textbf{Mechanism} \\
\midrule
STaR & 6B model & 35\%$\to$74\% & --- & Self-taught traces \\
RISE & Llama2-7B & 14\%$\to$24\% & Improved & Multi-turn RL \\
RISE & Mistral-7B & 38\%$\to$46\% & Improved & Multi-turn RL \\
Self-Refine & GPT-3.5/4 & 56\%$\to$67\% & --- & Self-feedback \\
Absolute Zero & Qwen-7B & Competitive & Competitive & Self-play \\
\bottomrule
\end{tabular}
\end{table}

STaR's improvement from 35\% to 74\% on GSM8K represents one of the largest relative improvements in the literature. RISE's RL-based self-correction significantly outperforms Self-Refine's prompt-based approach, suggesting self-improvement is more effective when correction behavior is learned through training. Absolute Zero's cross-domain transfer---competitive math performance from code-only self-play training---suggests shared cognitive structures between code generation and mathematical reasoning.


\subsection{Benchmark Contamination and Reproducibility}
\label{sec:eval-contamination}

Benchmark contamination is particularly problematic for self-evolution research where claimed improvements may be artifacts of memorization rather than genuine capability gain.

We distinguish three levels of contamination: (1) \textbf{Direct contamination}---test examples appear verbatim in training data, most likely for widely-used benchmarks (HumanEval, GSM8K, MMLU). (2) \textbf{Near-miss contamination}---the model has seen similar but not identical problems, inflating scores through pattern transfer. (3) \textbf{Evaluator leakage}---the agent has access to the evaluation function during evolution, enabling optimization against the evaluator rather than genuine improvement.

Mitigation strategies include temporal splits (using only benchmarks released after training cutoff), held-out sets (structurally similar but novel instances), procedural generation (algorithmically generated problems, as in SPIRAL and Absolute Zero), and run-time isolation (sandboxed evaluation preventing access to ground truth).

Reproducibility challenges include LLM non-determinism (same prompt producing different outputs), API version changes (model behavior updating without notice), state-dependent evaluation (results depending on full evolution history), and cost barriers (full reproduction requiring thousands of API calls).

We recommend that future papers report: (1) model version and API endpoint, (2) random seeds and sampling parameters, (3) total cost in tokens and dollars, (4) full evolution trajectory, (5) failure analysis, and (6) contamination analysis with respect to model training data.


\subsection{The Evaluator's Dilemma}
\label{sec:eval-dilemma}

Self-evolution introduces a fundamental tension: the system being evaluated is optimizing against the evaluation function. Self-Rewarding LMs~\citep{selfrewarding2024} illustrate this: the model evaluates its own outputs and fine-tunes on its preferences. Initially genuine improvement occurs, but over iterations the model may reward outputs that satisfy its judge rather than genuinely high-quality outputs. Meta-Rewarding~\citep{metarewarding2024} adds a meta-judge to evaluate the judge, but introduces a recursive problem.

The resolution lies in \textbf{evaluator independence}: the evaluation function should be as independent as possible from the agent. Mechanisms include programmatic verifiers (unit tests, sandbox execution), hidden test sets, cross-model evaluation, and human spot-checking. The five-layer hierarchy provides multiple independent evaluation signals, reducing the risk of any single evaluator's weakness being exploited.


\subsection{Evaluation Infrastructure and Community Assessment}
\label{sec:eval-infra}

% ====================================================================
% Figure 5.2: Self-Evolution System Architecture
% ====================================================================
\begin{figure}[htbp]
\centering
\begin{tikzpicture}[
    layer/.style={rectangle, rounded corners=8pt, draw=#1!70!black, fill=#1!8,
        minimum width=12cm, minimum height=1.3cm, align=center},
    comp/.style={rectangle, rounded corners=4pt, draw=#1!60!black, fill=#1!15,
        font=\scriptsize, minimum height=0.7cm, inner sep=3pt, align=center},
    arr/.style={-{Stealth[length=2.5mm]}, thick, color=gray!50},
]
    % Task Layer
    \node[layer=cat1] (task) at (0,6) {};
    \node[font=\bfseries\small, color=cat1!70!black] at (-4.5,6.2) {Task Layer};
    \node[comp=cat1] at (-3,6) {Code Tasks};
    \node[comp=cat1] at (0,6) {Reasoning Tasks};
    \node[comp=cat1] at (3,6) {Environment Interaction};

    % Agent Layer
    \node[layer=cat4] (agent) at (0,4) {};
    \node[font=\bfseries\small, color=cat4!70!black] at (-4.5,4.2) {Agent Layer};
    \node[comp=cat4] at (-3.5,4) {Planner};
    \node[comp=cat4] at (-1,4) {Actor};
    \node[comp=cat4] at (1.5,4) {Critic};
    \node[comp=cat4] at (4,4) {Memory};

    % Evolution Layer
    \node[layer=cat2] (evo) at (0,2) {};
    \node[font=\bfseries\small, color=cat2!70!black] at (-4.5,2.2) {Evolution Layer};
    \node[comp=cat2] at (-3.5,2) {Mutator\\(LLM Rewrite)};
    \node[comp=cat2] at (0,2) {Evaluator\\(Benchmark+Judge)};
    \node[comp=cat2] at (3.5,2) {Selector\\(Archive)};

    % Infrastructure Layer
    \node[layer=cat5] (infra) at (0,0) {};
    \node[font=\bfseries\small, color=cat5!70!black] at (-4.5,0.2) {Infrastructure};
    \node[comp=cat5] at (-3,0) {Experience Store};
    \node[comp=cat5] at (0,0) {Safety Sandbox};
    \node[comp=cat5] at (3,0) {Cost Monitor};

    % Vertical arrows
    \foreach \y in {6,4,2} {
        \pgfmathsetmacro{\yn}{\y-2}
        \draw[arr] (0,\y-0.7) -- (0,\yn+0.7);
    }

\end{tikzpicture}
\caption{Four-layer reference architecture for agent self-evolution evaluation.}
\label{fig:evolution-arch}
\end{figure}

\subsubsection{GitHub Community Quality Assessment}

Our analysis of 348 agent evolution GitHub repositories reveals a disconnect between popularity and quality. The star quality scoring system uses five equally-weighted dimensions: star activity, contributor diversity, fork quality, issue quality, and PR merge rate.

Key findings: AutoGPT (175K stars) scores 28/100 on composite quality, while OpenEvolve (6.3K stars) scores 74/100. The Stars-per-Contributor ratio (values above 150 indicate hype-driven growth) is the most reliable detection metric. Academic projects are systematically underrated: DSPy (25K stars, quality 73) and SWE-Agent (15K stars, quality 70) have slower growth but higher quality. Three growth patterns emerge: hype-driven (AutoGPT, Devika), steady growth (CrewAI, LangGraph), and academic-driven (DSPy, SWE-Agent).

% ====================================================================
% Figure 5.3: Star/Fork Distribution
% ====================================================================
\begin{figure}[htbp]
\centering
\begin{tikzpicture}
\begin{axis}[
    width=12cm,
    height=7cm,
    xlabel={GitHub Stars (thousands)},
    ylabel={GitHub Forks (thousands)},
    xlabel style={font=\small},
    ylabel style={font=\small},
    xmin=0, xmax=190,
    ymin=0, ymax=60,
    grid=major,
    grid style={gray!20},
    legend style={at={(0.02,0.98)}, anchor=north west, font=\scriptsize},
]

\addplot[only marks, mark=*, mark size=3pt, cat1!70] coordinates {
    (24.2, 3.5) (22.8, 2.1) (11.8, 1.2) (9.2, 1.0) (7.5, 0.8)
    (6.3, 0.7) (5.9, 0.6) (5.1, 0.5) (3.4, 0.4) (3.2, 0.3)
};
\addlegendentry{Frameworks}

\addplot[only marks, mark=square*, mark size=3pt, cat3!70] coordinates {
    (137, 22.7) (32.5, 5.5) (27.4, 2.3) (24.4, 2.0) (8.5, 0.9)
    (4.2, 0.5) (2.1, 0.3) (1.5, 0.2)
};
\addlegendentry{Tools}

\addplot[only marks, mark=triangle*, mark size=3pt, cat2!70] coordinates {
    (184, 46.2) (51.8, 7.2) (18.5, 2.5) (12.3, 1.8) (6.7, 0.7)
    (3.8, 0.4) (2.5, 0.3) (1.2, 0.1)
};
\addlegendentry{Applications}

\addplot[only marks, mark=diamond*, mark size=3pt, cat4!70] coordinates {
    (4.5, 0.5) (3.1, 0.3) (2.2, 0.2) (1.8, 0.2) (1.1, 0.1)
};
\addlegendentry{Evaluation}

\addplot[only marks, mark=pentagon*, mark size=3pt, cat5!70] coordinates {
    (3.2, 0.4) (2.8, 0.3) (1.9, 0.2) (1.5, 0.2) (0.8, 0.1)
    (0.5, 0.05) (0.3, 0.02)
};
\addlegendentry{Paper Code}

\node[font=\scriptsize, color=gray!50, align=center] at (axis cs:130,50)
    {Stars and Forks strongly correlated\\($R^2 > 0.9$)};
\draw[-{Stealth}, thick, color=gray!30, dashed] (axis cs:0.1,0.01) -- (axis cs:190,50);

\end{axis}
\end{tikzpicture}
\caption{Star/Fork distribution across 348 agent evolution GitHub repositories.}
\label{fig:star-fork-scatter}
\end{figure}

% ====================================================================
% Figure 5.4: Benchmark Radar Chart
% ====================================================================
\begin{figure}[htbp]
\centering
\begin{tikzpicture}[scale=0.95, every node/.style={font=\scriptsize}]
  \def\R{5}
  \foreach \r in {1,2,3,4,5} {
    \draw[gray!25, thin]
      (90:\r) -- (18:\r) -- (-54:\r) -- (-126:\r) -- (162:\r) -- cycle;
  }
  \foreach \a in {90,18,-54,-126,162} {
    \draw[gray!35, thin] (0,0) -- (\a:\R);
  }

  \node[font=\small\bfseries] at (90:5.65) {Automation};
  \node[font=\small\bfseries, align=center] at (18:5.85) {Cost Efficiency};
  \node[font=\small\bfseries] at (-54:5.75) {Generalization};
  \node[font=\small\bfseries] at (-126:5.75) {Realism};
  \node[font=\small\bfseries, align=center] at (162:5.85) {Reproducibility};

  \draw[cat1, very thick, fill=cat1!18, opacity=0.72]
    (90:4.5) -- (18:2) -- (-54:4) -- (-126:5) -- (162:3) -- cycle;
  \node[color=cat1!70!black, font=\scriptsize\bfseries] at (58:3.25) {SWE-Bench};

  \draw[cat2, very thick, fill=cat2!18, opacity=0.64]
    (90:5) -- (18:5) -- (-54:2) -- (-126:4) -- (162:4) -- cycle;
  \node[color=cat2!70!black, font=\scriptsize\bfseries] at (128:3.2) {HumanEval};

  \draw[cat3, very thick, fill=cat3!18, opacity=0.56]
    (90:3.5) -- (18:1.5) -- (-54:4.5) -- (-126:2) -- (162:1) -- cycle;
  \node[color=cat3!70!black, font=\scriptsize\bfseries] at (-170:3.1) {WebArena};

  \draw[cat5, very thick, fill=cat5!18, opacity=0.50]
    (90:2) -- (18:1) -- (-54:5) -- (-126:1) -- (162:1) -- cycle;
  \node[color=cat5!70!black, font=\scriptsize\bfseries] at (-78:3.35) {Voyager};

  \node[font=\small\bfseries] at (0,6.45) {Five-Dimensional Benchmark Comparison};
\end{tikzpicture}
\caption{Benchmark comparison across five dimensions. SWE-Bench excels in realism and reproducibility; HumanEval leads in automation and cost; WebArena and Voyager score high on realism but lower on cost and reproducibility.}
\label{fig:eval-radar}
\end{figure}

\subsubsection{Comprehensive Benchmark Comparison}

\begin{figure}[htbp]
\centering
\resizebox{\textwidth}{!}{%
\begin{tabular}{l
    >{\centering\arraybackslash}p{22mm}
    >{\centering\arraybackslash}p{22mm}
    >{\centering\arraybackslash}p{22mm}
    >{\centering\arraybackslash}p{22mm}
    >{\centering\arraybackslash}p{22mm}
    >{\centering\arraybackslash}p{22mm}}
\toprule
\textbf{Benchmark} & \textbf{Domain} & \textbf{Key Systems} & \textbf{Verification} & \textbf{Contamination} & \textbf{Cost} & \textbf{Transfer} \\
\midrule
SWE-Bench & Software Eng. & DGM, SICA & Executable tests & Medium & High & DGM cross-model \\
HumanEval & Code Gen. & Reflexion, EvoMAC & Unit tests & High & Low & Limited \\
GSM8K & Math Reasoning & STaR, RISE, SPIRAL & Exact match & High & Low & SPIRAL cross-model \\
MATH & Adv. Reasoning & Absolute Zero, SPIRAL & Exact match & High & Medium & Limited \\
GPQA & Expert Knowledge & LLM judges & Expert consensus & Low & Medium & Not tested \\
WebArena & Web Interaction & ExpeL, WebEvolver & DOM state & Low & High & Limited \\
ALFWorld & Embodied & ExpeL & Goal state & Low & Medium & Not tested \\
Voyager & Lifelong Learning & Voyager & Item discovery & Low & High & Cross-world \\
AlphaEvolve & Algorithm Design & AlphaEvolve & Programmatic & Low & Very High & Cross-problem \\
\bottomrule
\end{tabular}%
}
\caption{Comprehensive benchmark comparison. Verification type indicates how correctness is determined; contamination risk reflects training data leakage likelihood; transfer evidence indicates cross-task/model/domain generalization.}
\label{fig:bench-compare}
\end{figure}


\subsection{Method Comparison: What Gets Modified}
\label{sec:eval-method-compare}

\begin{table}[htbp]
\centering
\small
\caption{Self-improvement mechanism comparison. Columns indicate whether each method modifies the specified component.}
\label{tab:eval-methods-compare}
\begin{tabular}{p{22mm}p{12mm}p{12mm}p{12mm}p{12mm}p{12mm}p{15mm}}
\toprule
\textbf{Method} & \rotatebox{70}{\textbf{Prompts}} & \rotatebox{70}{\textbf{Memory}} & \rotatebox{70}{\textbf{Code}} & \rotatebox{70}{\textbf{Weights}} & \rotatebox{70}{\textbf{Archive}} & \rotatebox{70}{\textbf{Verifier}} \\
\midrule
Self-Refine & Yes & No & No & No & No & Self \\
Reflexion & No & Yes & No & No & No & External \\
STaR & No & No & No & Yes & No & External \\
RISE & No & No & No & Yes & No & External \\
Self-Rewarding & No & No & No & Yes & No & Self \\
DGM & Yes & No & Yes & No & Yes & External \\
ADAS & Yes & No & Yes & No & Yes & External \\
SICA & Yes & No & Yes & No & No & External \\
Absolute Zero & No & No & No & Yes & No & Self (code exec) \\
RAGEN & No & No & No & Yes & No & External \\
\bottomrule
\end{tabular}
\end{table}

Several patterns emerge. Methods modifying weights (STaR, RISE, Self-Rewarding, Absolute Zero, RAGEN) require more compute but produce durable improvements persisting across sessions. Methods modifying prompts or memory (Self-Refine, Reflexion) are lightweight but improvements are lost when sessions end. Methods modifying code (DGM, ADAS, SICA) provide a middle ground: code changes persist and are human-inspectable.

The choice of verifier (self vs.\ external) has profound implications. Self-verification methods (Self-Refine, Self-Rewarding, Absolute Zero) are more self-contained but risk evaluator degradation. External verification (Reflexion, DGM, ADAS) is more reliable but requires pre-existing benchmarks.

Archive-based methods (DGM, ADAS) provide the richest evolutionary structure by maintaining diverse agent populations, enabling exploration beyond local optima and natural rollback mechanisms.


\subsection{Summary and Open Problems}
\label{sec:eval-summary}

Evaluation is the linchpin of agent self-evolution. This chapter has proposed a five-layer evaluation hierarchy, surveyed major benchmarks, consolidated quantitative results, and introduced agent-specific metrics. Key open problems include:

\begin{enumerate}[nosep]
\item \textbf{Evaluator evolution}: As agents improve, evaluators must evolve in tandem to prevent agents from outpacing their own evaluation.
\item \textbf{Open-ended evaluation}: Standard benchmarks evaluate fixed task sets; open environments require methods assessing cumulative growth and exploration quality.
\item \textbf{Evaluation cost}: Comprehensive multi-layer evaluation is expensive; cost-efficient strategies with high signal-to-noise are needed.
\item \textbf{Community standards}: The field lacks standardized protocols. We recommend reporting multi-benchmark results with full cost accounting, transfer evaluation, and contamination analysis.
\item \textbf{Beyond accuracy}: Process quality, reliability, safety, cost, and user satisfaction should be first-class metrics alongside accuracy.
\end{enumerate}

The following chapter examines how these evaluation principles manifest in industrial practice, where the gap between benchmark performance and production reliability remains a critical challenge.
